% ===========================================================================
%  Programmation Frontend — Lab 2 : Composants, Props et CSS Modules
%  ESP/UCAD — 2025–2026
%  Dr. El Hadji Bassirou TOURÉ
%  Compilé avec : xelatex (2–3 passes)
% ===========================================================================
\documentclass[11pt,a4paper]{article}

\usepackage{fontspec}
\setmainfont{Latin Modern Roman}
\setsansfont{Latin Modern Sans}
\setmonofont{Latin Modern Mono}
\usepackage{polyglossia}
\setdefaultlanguage{french}

\usepackage[margin=2.5cm]{geometry}
\usepackage{amsmath,amssymb}
\usepackage{graphicx}
\usepackage{xcolor}
\usepackage{booktabs}
\usepackage{hyperref}
\usepackage{tcolorbox}
\tcbuselibrary{theorems,breakable,skins}
\usepackage{enumitem}
\usepackage{fancyhdr}
\usepackage{array}
\usepackage{pifont}
\usepackage{longtable}
\usepackage{multirow}

\newcommand{\cmark}{\ding{51}}
\newcommand{\xmark}{\ding{55}}

\hypersetup{
    colorlinks=true,
    linkcolor=blue!60!black,
    urlcolor=blue!60!black,
    citecolor=blue!60!black
}

\newtcolorbox{objectifbox}[1][]{colback=blue!5,colframe=blue!50,fonttitle=\bfseries,title=#1,breakable}
\newtcolorbox{infobox}[1][]{colback=green!5,colframe=green!50!black,fonttitle=\bfseries,title=#1,breakable}
\newtcolorbox{attentionbox}[1][]{colback=orange!5,colframe=orange!70,fonttitle=\bfseries,title=#1,breakable}
\newtcolorbox{retenirbox}[1][]{colback=yellow!10,colframe=yellow!50!black,fonttitle=\bfseries,title=#1,breakable}
\newtcolorbox{prereqbox}[1][]{colback=cyan!5,colframe=cyan!50!black,fonttitle=\bfseries,title=#1,breakable}
\newtcolorbox{projetbox}[1][]{colback=teal!5,colframe=teal!50!black,fonttitle=\bfseries,title=#1,breakable}

% --- Code source ---
\usepackage{listings}
\lstdefinelanguage{JavaScript}{
  keywords={break, case, catch, continue, debugger, default, delete, do, else,
    finally, for, function, if, in, instanceof, new, return, switch, this,
    throw, try, typeof, var, void, while, with, const, let, class, export,
    extends, import, super, yield, async, await, of, from, true, false, null,
    undefined},
  morecomment=[l]{//},
  morecomment=[s]{/*}{*/},
  morestring=[b]',
  morestring=[b]",
  morestring=[b]`,
  sensitive=true,
}

\lstdefinelanguage{CSS}{
  keywords={color, background, border, padding, margin, font, display, flex,
    align, justify, width, height, gap, cursor, none, solid, center, column,
    row, wrap, pointer, radius, size, weight, family, items, content,
    direction, transform, transition, hover, box, shadow, max, min},
  morecomment=[s]{/*}{*/},
  morestring=[b]',
  morestring=[b]",
  sensitive=true,
}

\lstset{
  basicstyle=\ttfamily\small,
  backgroundcolor=\color{gray!8},
  frame=single,
  framerule=0.5pt,
  rulecolor=\color{gray!40},
  breaklines=true,
  breakatwhitespace=true,
  showstringspaces=false,
  tabsize=2,
  xleftmargin=4pt,
  xrightmargin=4pt,
  aboveskip=8pt,
  belowskip=8pt,
  language=JavaScript,
  keywordstyle=\color{blue!70!black}\bfseries,
  commentstyle=\color{green!50!black}\itshape,
  stringstyle=\color{red!60!black},
  literate={é}{{\'e}}1 {è}{{\`e}}1 {ê}{{\^e}}1 {à}{{\`a}}1 {ù}{{\`u}}1
           {ô}{{\^o}}1 {î}{{\^i}}1 {ç}{{\c{c}}}1 {É}{{\'E}}1
           {→}{{$\rightarrow$}}1
}

\lstdefinestyle{terminal}{
  language={},
  keywordstyle={},
  commentstyle={},
  stringstyle={},
  basicstyle=\ttfamily\small,
  backgroundcolor=\color{black!5},
}

\lstdefinestyle{css}{
  language=CSS,
  keywordstyle=\color{blue!70!black}\bfseries,
  commentstyle=\color{green!50!black}\itshape,
  stringstyle=\color{red!60!black},
}

\pagestyle{fancy}
\fancyhf{}
\fancyhead[L]{\small Prog. Frontend --- Lab 2 : Composants, Props et CSS Modules}
\fancyhead[R]{\small ESP/UCAD}
\fancyfoot[C]{\thepage}
\fancyfoot[L]{\small Dr. El Hadji Bassirou TOURÉ}
\fancyfoot[R]{\small 2025--2026}
\renewcommand{\headrulewidth}{0.4pt}
\renewcommand{\footrulewidth}{0.4pt}

\title{%
\vspace{-1cm}
\textbf{\LARGE Programmation Frontend}\\[0.8cm]
\textbf{\huge Lab 2 --- Composants, Props et CSS Modules}\\[0.3cm]
\textbf{\Large Organiser et styliser une application React}\\[0.5cm]
\large ESP/UCAD --- 2025--2026}
\author{Dr. El Hadji Bassirou TOURÉ\\
Département Génie Informatique\\
Université Cheikh Anta Diop de Dakar}
\date{}

\begin{document}

\maketitle
\thispagestyle{empty}

\vspace{0.5cm}

% ===========================================================================
% BOÎTE PRÉREQUIS
% ===========================================================================
\begin{prereqbox}[Informations pratiques]
\begin{center}
\begin{tabular}{lp{10cm}}
\toprule
\textbf{Durée estimée} & 2h (une séance) \\
\textbf{Prérequis} & Avoir terminé le Lab 1 (dossier \texttt{annuaire-react/} fonctionnel). \\
\textbf{Fichier de départ} & Le fichier \texttt{src/App.jsx} du Lab 1 \\
\textbf{Livrables} & 3 composants (\texttt{App}, \texttt{UserCard}, \texttt{SearchBar}) avec CSS Modules \\
\textbf{Validation} & Le filtre en temps réel fonctionne, chaque composant a son propre fichier \\
\bottomrule
\end{tabular}
\end{center}
\end{prereqbox}

\begin{objectifbox}[Objectifs du lab]
À l'issue de ce lab, chaque étudiant sera capable de :
\begin{enumerate}[nosep]
    \item \textbf{Découper} une application React en composants réutilisables
    \item \textbf{Passer des données} d'un composant parent à un composant enfant via les \textbf{props}
    \item \textbf{Styliser} chaque composant avec un \textbf{CSS Module} (pas de collision de noms)
    \item \textbf{Remonter l'état} vers le parent commun (\emph{lifting state up})
    \item Implémenter un \textbf{filtre en temps réel} sur une liste
\end{enumerate}
\end{objectifbox}

\tableofcontents

\newpage

% ===========================================================================
\section{Point de départ --- Le code du Lab 1}
\label{sec:starter}
% ===========================================================================

Au Lab 1, vous avez créé une application React qui récupère des utilisateurs depuis JSONPlaceholder et les affiche en cartes. Tout fonctionne. Mais regardez le fichier \texttt{src/App.jsx} :

\begin{lstlisting}
import { useState } from "react";

const App = () => {
  const [users, setUsers] = useState([]);
  const [chargement, setChargement] = useState(false);

  const chargerUtilisateurs = async () => {
    setChargement(true);
    try {
      const response = await fetch(
        "https://jsonplaceholder.typicode.com/users"
      );
      const data = await response.json();
      setUsers(data);
    } catch (error) {
      console.error("Erreur :", error);
    }
    setChargement(false);
  };

  return (
    <div style={{ maxWidth: "600px", margin: "2rem auto",
                  fontFamily: "sans-serif" }}>
      <h1>Annuaire des utilisateurs</h1>
      <button onClick={chargerUtilisateurs}
              disabled={chargement}>
        {chargement ? "Chargement..." : "Charger les utilisateurs"}
      </button>

      {users.map(user => (
        <div key={user.id} style={{ border: "1px solid #ccc",
             padding: "1rem", margin: "0.5rem 0",
             borderRadius: "8px" }}>
          <h3 style={{ margin: "0 0 0.5rem 0" }}>{user.name}</h3>
          <p style={{ margin: "0.2rem 0", color: "#555" }}>
            Email : {user.email}
          </p>
          <p style={{ margin: "0.2rem 0", color: "#555" }}>
            Entreprise : {user.company.name}
          </p>
        </div>
      ))}
    </div>
  );
};

export default App;
\end{lstlisting}

\begin{objectifbox}[Étape 0 --- Vérification]
Ouvrez le dossier \texttt{annuaire-react/} dans VS Code. Lancez le serveur de développement :

\begin{lstlisting}[style=terminal]
cd annuaire-react
npm run dev
\end{lstlisting}

Ouvrez \texttt{http://localhost:5173}. Cliquez sur le bouton. Les 10 utilisateurs doivent apparaître en cartes.

Si le projet ne démarre pas, vérifiez que les dépendances sont installées (\texttt{npm install}).
\end{objectifbox}

\subsection{Le problème}

Le code fonctionne, mais tout est dans un seul fichier :

\begin{itemize}[nosep]
    \item La logique de chargement (fetch, loading)
    \item L'affichage de chaque carte utilisateur (le bloc \texttt{<div>} dans le \texttt{.map()})
    \item Les styles (en \texttt{style=\{\{...\}\}} directement dans le JSX)
\end{itemize}

Imaginez que Fatou doit ajouter un champ de recherche, que Moussa doit modifier l'apparence des cartes, et que Cheikh doit ajouter un compteur. Les trois travaillent dans le \textbf{même fichier}. Ça ne tient pas.

La solution : \textbf{découper en composants}. Chaque morceau d'interface devient un fichier indépendant avec sa propre logique et son propre style.

% ===========================================================================
\section{Partie 1 --- Extraire un composant}
\label{sec:composant}
% ===========================================================================

\begin{objectifbox}[Objectif --- 15 minutes]
Extraire la carte utilisateur dans un composant \texttt{UserCard} séparé.
\end{objectifbox}

\subsection{Créer le dossier \texttt{components/}}

Dans \texttt{src/}, créez un dossier \texttt{components/} :

\begin{lstlisting}[style=terminal]
mkdir src/components
\end{lstlisting}

Ce dossier contiendra tous les composants réutilisables de l'application.

\subsection{Créer \texttt{UserCard.jsx}}

Créez le fichier \texttt{src/components/UserCard.jsx} :

\begin{lstlisting}
const UserCard = ({ user }) => {
  return (
    <div style={{ border: "1px solid #ccc", padding: "1rem",
                  margin: "0.5rem 0", borderRadius: "8px" }}>
      <h3 style={{ margin: "0 0 0.5rem 0" }}>{user.name}</h3>
      <p style={{ margin: "0.2rem 0", color: "#555" }}>
        Email : {user.email}
      </p>
      <p style={{ margin: "0.2rem 0", color: "#555" }}>
        Entreprise : {user.company.name}
      </p>
    </div>
  );
};

export default UserCard;
\end{lstlisting}

On a simplement \textbf{coupé} le bloc \texttt{<div>} qui était dans le \texttt{.map()} et on l'a mis dans sa propre fonction. Le composant reçoit un objet \texttt{user} et affiche ses informations.

\subsection{Utiliser \texttt{UserCard} dans \texttt{App.jsx}}

Modifiez \texttt{src/App.jsx} pour importer et utiliser le nouveau composant :

\begin{lstlisting}
import { useState } from "react";
import UserCard from "./components/UserCard";

const App = () => {
  const [users, setUsers] = useState([]);
  const [chargement, setChargement] = useState(false);

  const chargerUtilisateurs = async () => {
    setChargement(true);
    try {
      const response = await fetch(
        "https://jsonplaceholder.typicode.com/users"
      );
      const data = await response.json();
      setUsers(data);
    } catch (error) {
      console.error("Erreur :", error);
    }
    setChargement(false);
  };

  return (
    <div style={{ maxWidth: "600px", margin: "2rem auto",
                  fontFamily: "sans-serif" }}>
      <h1>Annuaire des utilisateurs</h1>
      <button onClick={chargerUtilisateurs}
              disabled={chargement}>
        {chargement ? "Chargement..." : "Charger les utilisateurs"}
      </button>

      {users.map(user => (
        <UserCard key={user.id} user={user} />
      ))}
    </div>
  );
};

export default App;
\end{lstlisting}

Sauvegardez les deux fichiers. La page se met à jour automatiquement. Le résultat est \textbf{identique} à avant, mais le code est mieux organisé.

\begin{attentionbox}[Vérification]
Si vous voyez une erreur du type \texttt{UserCard is not defined} : vérifiez que l'import est correct. Le chemin est relatif : \texttt{"./components/UserCard"} (avec le \texttt{./} au début).
\end{attentionbox}

% ===========================================================================
\section{Partie 2 --- Les props}
\label{sec:props}
% ===========================================================================

\begin{objectifbox}[Objectif --- 15 minutes]
Comprendre comment les données circulent entre composants grâce aux props.
\end{objectifbox}

\subsection{Comment les données circulent}

Dans le code qu'on vient d'écrire, regardez cette ligne dans \texttt{App.jsx} :

\begin{lstlisting}
<UserCard key={user.id} user={user} />
\end{lstlisting}

Et cette ligne dans \texttt{UserCard.jsx} :

\begin{lstlisting}
const UserCard = ({ user }) => {
\end{lstlisting}

Le mot \texttt{user} apparaît des deux côtés. Voici ce qui se passe :

\begin{enumerate}[nosep]
    \item \texttt{App} \textbf{envoie} l'objet \texttt{user} au composant \texttt{UserCard} via l'attribut \texttt{user=\{user\}}.
    \item \texttt{UserCard} \textbf{reçoit} cet objet dans ses \textbf{props} (propriétés).
    \item L'écriture \texttt{(\{ user \})} est une \textbf{déstructuration} : on extrait directement la propriété \texttt{user} de l'objet props.
\end{enumerate}

\begin{infobox}[Analogie : la fiche de commande]
Imaginez un restaurant. Le serveur (le composant parent) remplit une fiche de commande et la passe au cuisinier (le composant enfant). La fiche contient les informations nécessaires : «~1 thiéboudienne, pas trop pimenté~». Le cuisinier lit la fiche et prépare le plat. Il ne modifie \textbf{pas} la fiche — il la lit seulement.

Les props fonctionnent pareil : le parent écrit, l'enfant lit. \textbf{Les props sont en lecture seule.}
\end{infobox}

\subsection{Props multiples}

Un composant peut recevoir autant de props que nécessaire. Modifions \texttt{UserCard} pour recevoir une prop supplémentaire qui contrôle l'affichage de l'email :

\begin{lstlisting}
const UserCard = ({ user, afficherEmail }) => {
  return (
    <div style={{ border: "1px solid #ccc", padding: "1rem",
                  margin: "0.5rem 0", borderRadius: "8px" }}>
      <h3 style={{ margin: "0 0 0.5rem 0" }}>{user.name}</h3>
      {afficherEmail && (
        <p style={{ margin: "0.2rem 0", color: "#555" }}>
          Email : {user.email}
        </p>
      )}
      <p style={{ margin: "0.2rem 0", color: "#555" }}>
        Entreprise : {user.company.name}
      </p>
    </div>
  );
};

export default UserCard;
\end{lstlisting}

Dans \texttt{App.jsx} :

\begin{lstlisting}
<UserCard key={user.id} user={user} afficherEmail={true} />
\end{lstlisting}

La syntaxe \texttt{\{afficherEmail \&\& (...)\}} signifie : «~si \texttt{afficherEmail} est vrai, affiche ce bloc~». C'est du \textbf{rendu conditionnel}.

\begin{retenirbox}[Les props en résumé]
\begin{center}
\begin{tabular}{lp{8cm}}
\toprule
\textbf{Concept} & \textbf{Explication} \\
\midrule
Passer une prop & \texttt{<Composant nom=\{valeur\} />} \\
Recevoir une prop & \texttt{const Composant = (\{ nom \}) => \{ ... \}} \\
Props multiples & \texttt{<Composant a=\{1\} b=\{"texte"\} c=\{true\} />} \\
Lecture seule & L'enfant ne modifie \textbf{jamais} les props reçues \\
\texttt{key} & Identifiant unique dans un \texttt{.map()} (pas une prop classique) \\
\bottomrule
\end{tabular}
\end{center}
\end{retenirbox}

Pour la suite du lab, retirez la prop \texttt{afficherEmail} et revenez à la version simple de \texttt{UserCard} (qui affiche toujours l'email). On ajoute cette prop uniquement pour illustrer le concept.

% ===========================================================================
\section{Partie 3 --- CSS Modules}
\label{sec:css-modules}
% ===========================================================================

\begin{objectifbox}[Objectif --- 15 minutes]
Remplacer les styles en ligne par des CSS Modules pour une meilleure organisation.
\end{objectifbox}

\subsection{Le problème des styles en ligne}

Dans le code actuel, les styles sont écrits directement dans le JSX :

\begin{lstlisting}
<div style={{ border: "1px solid #ccc", padding: "1rem",
              margin: "0.5rem 0", borderRadius: "8px" }}>
\end{lstlisting}

Trois problèmes :
\begin{itemize}[nosep]
    \item Le code est \textbf{difficile à lire} (mélange de HTML et de CSS).
    \item On ne peut pas utiliser les \textbf{pseudo-classes} (\texttt{:hover}, \texttt{:focus}).
    \item Si deux composants utilisent la même classe CSS, il y a un risque de \textbf{collision de noms}.
\end{itemize}

\subsection{CSS Modules : du CSS scopé par composant}

Un \textbf{CSS Module} est un fichier CSS classique dont les noms de classes sont \textbf{automatiquement rendus uniques} par Vite. Chaque composant a son propre fichier \texttt{.module.css}. Pas de collision possible.

Créez le fichier \texttt{src/components/UserCard.module.css} :

\begin{lstlisting}[style=css]
.card {
  border: 1px solid #ccc;
  padding: 1rem;
  margin: 0.5rem 0;
  border-radius: 8px;
  transition: box-shadow 0.2s;
}

.card:hover {
  box-shadow: 0 2px 8px rgba(0, 0, 0, 0.1);
}

.nom {
  margin: 0 0 0.5rem 0;
}

.info {
  margin: 0.2rem 0;
  color: #555;
}
\end{lstlisting}

Notez le \texttt{:hover} : on peut maintenant utiliser les pseudo-classes, ce qui était impossible avec les styles en ligne.

\subsection{Utiliser le CSS Module dans le composant}

Modifiez \texttt{src/components/UserCard.jsx} :

\begin{lstlisting}
import styles from "./UserCard.module.css";

const UserCard = ({ user }) => {
  return (
    <div className={styles.card}>
      <h3 className={styles.nom}>{user.name}</h3>
      <p className={styles.info}>Email : {user.email}</p>
      <p className={styles.info}>
        Entreprise : {user.company.name}
      </p>
    </div>
  );
};

export default UserCard;
\end{lstlisting}

\begin{infobox}[Comment fonctionne un CSS Module]
\begin{enumerate}[nosep]
    \item Vous écrivez \texttt{.card} dans le fichier CSS.
    \item Vite le transforme en quelque chose comme \texttt{.UserCard\_card\_x7k2q} (un nom unique).
    \item Vous y accédez avec \texttt{styles.card} dans le JSX.
    \item Résultat : même si un autre composant a aussi une classe \texttt{.card}, il n'y a \textbf{aucune collision}.
\end{enumerate}

Notez aussi : en React, on écrit \texttt{className} au lieu de \texttt{class} (car \texttt{class} est un mot réservé en JavaScript).
\end{infobox}

\subsection{CSS Module pour App}

Créez \texttt{src/App.module.css} :

\begin{lstlisting}[style=css]
.container {
  max-width: 600px;
  margin: 2rem auto;
  font-family: sans-serif;
  padding: 0 1rem;
}

.titre {
  color: #333;
}

.bouton {
  padding: 0.5rem 1.5rem;
  font-size: 1rem;
  cursor: pointer;
  border: 1px solid #ccc;
  border-radius: 6px;
  background: #f9f9f9;
  transition: background 0.2s;
}

.bouton:hover {
  background: #eee;
}

.bouton:disabled {
  cursor: not-allowed;
  opacity: 0.6;
}
\end{lstlisting}

Modifiez \texttt{src/App.jsx} :

\begin{lstlisting}
import { useState } from "react";
import UserCard from "./components/UserCard";
import styles from "./App.module.css";

const App = () => {
  const [users, setUsers] = useState([]);
  const [chargement, setChargement] = useState(false);

  const chargerUtilisateurs = async () => {
    setChargement(true);
    try {
      const response = await fetch(
        "https://jsonplaceholder.typicode.com/users"
      );
      const data = await response.json();
      setUsers(data);
    } catch (error) {
      console.error("Erreur :", error);
    }
    setChargement(false);
  };

  return (
    <div className={styles.container}>
      <h1 className={styles.titre}>Annuaire des utilisateurs</h1>
      <button className={styles.bouton}
              onClick={chargerUtilisateurs}
              disabled={chargement}>
        {chargement ? "Chargement..." : "Charger les utilisateurs"}
      </button>

      {users.map(user => (
        <UserCard key={user.id} user={user} />
      ))}
    </div>
  );
};

export default App;
\end{lstlisting}

Sauvegardez et vérifiez dans le navigateur. L'application fonctionne comme avant, avec en plus un effet \texttt{:hover} sur les cartes et sur le bouton.

\begin{retenirbox}[Avant / après : les styles]
\begin{center}
\begin{tabular}{lp{5cm}p{5cm}}
\toprule
& \textbf{Styles en ligne} (Lab 1) & \textbf{CSS Modules} (Lab 2) \\
\midrule
\textbf{Syntaxe} & \texttt{style=\{\{ color: "red" \}\}} & \texttt{className=\{styles.nom\}} \\
\textbf{Pseudo-classes} & Impossible & \texttt{:hover}, \texttt{:focus}, \texttt{:disabled} \\
\textbf{Collision} & Aucun risque (mais pas pratique) & Aucun risque (noms rendus uniques) \\
\textbf{Lisibilité} & Mélange JS + CSS & Séparation claire \\
\bottomrule
\end{tabular}
\end{center}
\end{retenirbox}

% ===========================================================================
\section{Partie 4 --- Créer le composant SearchBar}
\label{sec:searchbar}
% ===========================================================================

\begin{objectifbox}[Objectif --- 20 minutes]
Créer un composant de recherche et le connecter à la liste pour filtrer les utilisateurs en temps réel.
\end{objectifbox}

\subsection{Où placer la recherche ?}

On veut qu'en tapant dans un champ de recherche, la liste se filtre instantanément. Pour cela, il faut que :

\begin{enumerate}[nosep]
    \item Le composant \texttt{SearchBar} connaisse le texte tapé par l'utilisateur.
    \item Le composant \texttt{App} connaisse aussi ce texte pour filtrer la liste.
\end{enumerate}

Le texte de recherche est donc un \textbf{état partagé} entre \texttt{SearchBar} et \texttt{App}. En React, quand deux composants ont besoin du même état, on le place dans leur \textbf{parent commun}. C'est le principe du \emph{lifting state up} (remonter l'état).

\begin{infobox}[Analogie : le tableau de la classe]
Imaginez que Fatou et Moussa doivent consulter la même information. Au lieu que chacun ait sa propre copie (qui risque de diverger), on écrit l'information au tableau --- visible par tous. Le tableau, c'est le composant parent. L'information, c'est l'état.
\end{infobox}

\subsection{Créer \texttt{SearchBar.module.css}}

Créez \texttt{src/components/SearchBar.module.css} :

\begin{lstlisting}[style=css]
.container {
  margin: 1rem 0;
}

.input {
  width: 100%;
  padding: 0.6rem 1rem;
  font-size: 1rem;
  border: 1px solid #ccc;
  border-radius: 6px;
  box-sizing: border-box;
}

.input:focus {
  outline: none;
  border-color: #666;
}
\end{lstlisting}

\subsection{Créer \texttt{SearchBar.jsx}}

Créez \texttt{src/components/SearchBar.jsx} :

\begin{lstlisting}
import styles from "./SearchBar.module.css";

const SearchBar = ({ recherche, onRecherche }) => {
  return (
    <div className={styles.container}>
      <input
        className={styles.input}
        type="text"
        placeholder="Rechercher un utilisateur..."
        value={recherche}
        onChange={e => onRecherche(e.target.value)}
      />
    </div>
  );
};

export default SearchBar;
\end{lstlisting}

Le composant reçoit \textbf{deux props} :
\begin{itemize}[nosep]
    \item \texttt{recherche} : la valeur actuelle du champ (vient du parent).
    \item \texttt{onRecherche} : une \textbf{fonction} que le parent fournit. Quand l'utilisateur tape, on appelle cette fonction avec le nouveau texte.
\end{itemize}

\begin{attentionbox}[Passer une fonction comme prop]
En React, on passe souvent des fonctions comme props. Ici, \texttt{onRecherche} est une fonction fournie par \texttt{App}. Quand \texttt{SearchBar} appelle \texttt{onRecherche("Fatou")}, c'est en réalité \texttt{App} qui reçoit \texttt{"Fatou"} et met à jour son état. Le composant enfant \textbf{informe} le parent, sans modifier l'état directement.
\end{attentionbox}

\subsection{Connecter \texttt{SearchBar} dans \texttt{App.jsx}}

Modifiez \texttt{src/App.jsx} :

\begin{lstlisting}
import { useState } from "react";
import UserCard from "./components/UserCard";
import SearchBar from "./components/SearchBar";
import styles from "./App.module.css";

const App = () => {
  const [users, setUsers] = useState([]);
  const [chargement, setChargement] = useState(false);
  const [recherche, setRecherche] = useState("");

  const chargerUtilisateurs = async () => {
    setChargement(true);
    try {
      const response = await fetch(
        "https://jsonplaceholder.typicode.com/users"
      );
      const data = await response.json();
      setUsers(data);
    } catch (error) {
      console.error("Erreur :", error);
    }
    setChargement(false);
  };

  const usersFiltres = users.filter(user =>
    user.name.toLowerCase().includes(
      recherche.toLowerCase()
    )
  );

  return (
    <div className={styles.container}>
      <h1 className={styles.titre}>Annuaire des utilisateurs</h1>
      <button className={styles.bouton}
              onClick={chargerUtilisateurs}
              disabled={chargement}>
        {chargement ? "Chargement..." : "Charger les utilisateurs"}
      </button>

      <SearchBar
        recherche={recherche}
        onRecherche={setRecherche}
      />

      <p className={styles.compteur}>
        {usersFiltres.length} utilisateur(s) trouvé(s)
      </p>

      {usersFiltres.map(user => (
        <UserCard key={user.id} user={user} />
      ))}
    </div>
  );
};

export default App;
\end{lstlisting}

Ajoutez la classe \texttt{.compteur} dans \texttt{src/App.module.css} :

\begin{lstlisting}[style=css]
.compteur {
  color: #888;
  font-size: 0.9rem;
  margin: 0.5rem 0;
}
\end{lstlisting}

\subsection{Tester le filtre}

Sauvegardez tout. Chargez les utilisateurs, puis tapez dans le champ de recherche. La liste se filtre \textbf{en temps réel} à chaque lettre tapée. Le compteur indique le nombre de résultats.

\begin{retenirbox}[Flux de données du filtre]
Voici comment les données circulent dans l'application :

\begin{enumerate}[nosep]
    \item L'utilisateur tape «~lea~» dans \texttt{SearchBar}.
    \item \texttt{SearchBar} appelle \texttt{onRecherche("lea")}, qui est en réalité \texttt{setRecherche("lea")} dans \texttt{App}.
    \item L'état \texttt{recherche} dans \texttt{App} passe à \texttt{"lea"}.
    \item React re-rend \texttt{App}. Le filtre \texttt{users.filter(...)} ne garde que les utilisateurs dont le nom contient «~lea~».
    \item Le tableau filtré est passé au \texttt{.map()}, qui ne crée des \texttt{UserCard} que pour les résultats.
\end{enumerate}

Le flux est \textbf{unidirectionnel} : les données descendent (parent $\to$ enfant via props), les événements remontent (enfant $\to$ parent via fonctions callback).
\end{retenirbox}

% ===========================================================================
\section{Bilan --- Structure du projet}
\label{sec:bilan}
% ===========================================================================

Votre projet a maintenant cette structure :

\begin{lstlisting}[style=terminal]
annuaire-react/
  src/
    App.jsx
    App.module.css
    components/
      UserCard.jsx
      UserCard.module.css
      SearchBar.jsx
      SearchBar.module.css
    main.jsx
  index.html
  package.json
\end{lstlisting}

Chaque composant a \textbf{deux fichiers} : un \texttt{.jsx} pour la logique et un \texttt{.module.css} pour le style. Si Fatou veut modifier l'apparence des cartes, elle ne touche que \texttt{UserCard.module.css}. Si Moussa veut changer le comportement de la recherche, il ne touche que \texttt{SearchBar.jsx}. Personne ne se marche dessus.

\begin{retenirbox}[Avant / après : la structure]
\begin{center}
\begin{tabular}{lp{5cm}p{5cm}}
\toprule
& \textbf{Lab 1} & \textbf{Lab 2} \\
\midrule
\textbf{Fichiers} & 1 fichier (\texttt{App.jsx}) & 6 fichiers (3 composants $\times$ 2) \\
\textbf{Styles} & En ligne (\texttt{style=\{\{...\}\}}) & CSS Modules (\texttt{.module.css}) \\
\textbf{Recherche} & Aucune & Filtre en temps réel \\
\textbf{Réutilisabilité} & Aucune & \texttt{UserCard} réutilisable partout \\
\textbf{Collaboration} & Difficile (1 fichier) & Facile (fichiers séparés) \\
\bottomrule
\end{tabular}
\end{center}
\end{retenirbox}

% ===========================================================================
\section{Aide-mémoire}
\label{sec:aide-memoire}
% ===========================================================================

\begin{center}
\begin{tabular}{lp{8cm}}
\toprule
\textbf{Concept} & \textbf{Syntaxe / Exemple} \\
\midrule
Créer un composant & \texttt{const MonComposant = () => <div>...</div>;} \\
Exporter & \texttt{export default MonComposant;} \\
Importer & \texttt{import MonComposant from "./MonComposant";} \\
\midrule
Passer une prop & \texttt{<Enfant nom=\{valeur\} />} \\
Recevoir une prop & \texttt{const Enfant = (\{ nom \}) => ...} \\
Prop fonction & \texttt{<Enfant onAction=\{maFonction\} />} \\
Rendu conditionnel & \texttt{\{condition \&\& <Element />\}} \\
\midrule
CSS Module (fichier) & \texttt{MonComposant.module.css} \\
CSS Module (import) & \texttt{import styles from "./MonComposant.module.css";} \\
CSS Module (usage) & \texttt{className=\{styles.maClasse\}} \\
\midrule
Filtrer un tableau & \texttt{tableau.filter(item => condition)} \\
Lifting state up & L'état partagé vit dans le parent commun \\
\bottomrule
\end{tabular}
\end{center}

\newpage

% ===========================================================================
\section{Résumé et conclusion}
% ===========================================================================

Dans ce lab, vous avez :

\begin{enumerate}[nosep]
    \item \textbf{Extrait} la carte utilisateur dans un composant \texttt{UserCard} réutilisable.
    \item Compris les \textbf{props} : des données en lecture seule passées du parent vers l'enfant.
    \item Remplacé les styles en ligne par des \textbf{CSS Modules} (un fichier \texttt{.module.css} par composant).
    \item Créé un composant \texttt{SearchBar} et implémenté un \textbf{filtre en temps réel}.
    \item Appliqué le \emph{lifting state up} : l'état de recherche vit dans \texttt{App} et descend vers \texttt{SearchBar} et la liste filtrée.
\end{enumerate}

Voici où vous en êtes dans la progression du cours :

\begin{center}
\begin{tabular}{lcp{7.5cm}}
\toprule
\textbf{Étape} & \textbf{Statut} & \textbf{Ce que vous savez faire} \\
\midrule
\textbf{DOM et événements} & \cmark{} & Sélecteurs, événements, modification du DOM \\
\textbf{Fonctions JS} & \cmark{} & Déclaration, expression, fléchée, \emph{hoisting} \\
\textbf{Asynchrone} & \cmark{} & \texttt{fetch}, \texttt{async/await}, gestion d'erreurs \\
\textbf{React --- les bases} & \cmark{} & Composant, JSX, \texttt{useState}, \texttt{.map()} \\
\textbf{Composants et props} & \cmark{} & Découpage, props, CSS Modules, lifting state up \\
\textbf{useEffect et cycle de vie} & $\to$ Lab 3 & Charger au montage, loading, détail \\
\textbf{Formulaires et routing} & $\to$ Lab 4 & Inputs contrôlés, React Router \\
\bottomrule
\end{tabular}
\end{center}

\begin{projetbox}[Et ensuite ?]
L'annuaire affiche des utilisateurs et permet de les filtrer. Mais il faut encore cliquer sur un bouton pour charger les données. Dans le \textbf{Lab 3}, vous apprendrez \texttt{useEffect} pour charger les données \textbf{automatiquement} au démarrage, gérer les états de chargement et d'erreur proprement, et créer une page de détail pour chaque utilisateur.
\end{projetbox}

% ===========================================================================
\section*{Exercice bonus --- Pour aller plus loin}
% ===========================================================================

\begin{infobox}[Bonus : un composant Header]
Si vous avez terminé en avance, créez un composant \texttt{Header.jsx} + \texttt{Header.module.css} qui contient le titre de l'application et le bouton de chargement. Indices :

\begin{enumerate}[nosep]
    \item Créez \texttt{src/components/Header.jsx} et \texttt{src/components/Header.module.css}.
    \item Le composant \texttt{Header} reçoit trois props : \texttt{onCharger} (la fonction de chargement), \texttt{chargement} (booléen), et \texttt{nbResultats} (nombre).
    \item Dans \texttt{App.jsx}, remplacez le \texttt{<h1>} et le \texttt{<button>} par \texttt{<Header />}.
    \item L'application doit fonctionner exactement comme avant.
\end{enumerate}

Cet exercice renforce le découpage en composants et les props (y compris passer des fonctions).
\end{infobox}

\vspace{1cm}

\begin{center}
\fbox{
\begin{minipage}{0.85\textwidth}
\centering
\vspace{0.5cm}
\textbf{\Large Lab 2 --- Composants, Props et CSS Modules --- Terminé \cmark}\\[0.3cm]
\normalsize
Vous savez découper une application React en composants,\\
passer des données avec les props et styliser avec CSS Modules.\\[0.2cm]
Prochaine étape : \texttt{useEffect}, chargement automatique et page de détail.\\[0.3cm]
\textit{Un composant = une responsabilité. Un fichier = un composant.}\\[0.5cm]
\end{minipage}
}
\end{center}

\end{document}
